\section{Evaluation}

We organize the evaluation around four program-level research questions---RQ1 (audited Bugs4Q defects), RQ2 (balanced injected positive control), RQ3 (false positives), and RQ4 (sensitivity floor)---plus a circuit-level-only diagnostic run (\S\ref{sec:diagnostic}). Each claim is tied to a dataset, baseline, and metric, and each dataset is introduced with its research question below. MorphQ is the only matched baseline; QMutPy is not evaluated (it assumes an existing test suite), and stack/assertion tools are related-only comparators (Table~\ref{tab:related-comparison}). The primary evidence is program-level; the circuit-level run is retained only as a diagnostic and is not used as a main claim.

\subsection{Execution and Evidence Pipeline}

All runs use Qiskit Aer under a pinned environment, a fixed master seed, a default shot budget for detection, and a per-pair wall-clock timeout. MorphQ is run on the same subjects under identical backends, shot budgets, seeds, and timeouts. Every headline number resolves to one row of a canonical CSV that is derived by \texttt{derive\_metrics.py} from run-isolated summary shards under a frozen run directory; the derivation enforces expected denominators, rejects missing or stale-mixed rows, and re-derives the canonical CSV to a byte-identical copy (\emph{canonical integrity check: PASS}). Table~\ref{tab:canonical-results} is the canonical view.

\begin{table*}[t]
\centering
\caption{Canonical results, regenerated from the frozen run's \texttt{canonical\_results.csv}. Intervals are Wilson 95\% CIs. MorphQ's program-level detections are $0$ by the testing level (Sec.~\ref{sec:relations}). $^{\dagger}$Coverage-by-construction (admit-if-detected), not a detection rate (Sec.~\ref{sec:threats}). $^{\S}$The 3 mined subjects (\S\ref{sec:rq1}) meet the same checks but lie outside the frozen run. $^{\flat}$Degenerate cell (matched seeds give $p\equiv1$); the noisy/calibrated rows are informative.}
\label{tab:canonical-results}
\scriptsize
\setlength{\tabcolsep}{2.4pt}
\begin{tabular}{@{}p{0.16\textwidth}p{0.20\textwidth}p{0.12\textwidth}p{0.13\textwidth}p{0.08\textwidth}p{0.16\textwidth}@{}}
\toprule
\textbf{Cohort} & \textbf{Benchmark} & \textbf{Method} & \textbf{Metric} & \textbf{Count} & \textbf{Rate [95\% CI]} \\
\midrule
Audited Bugs4Q & Program-level defect subset & QMTester & Detected & 7/7 & 1.000 [0.646, 1.000] \\
Real defects (all)$^{\S}$ & Bugs4Q $+$ 3 mined (QCSE) & QMTester & Detected & \textbf{10/10} & \textbf{1.000 [0.722, 1.000]} \\
Audited Bugs4Q & Program-level defect subset & MorphQ & Detected & 0/7 & 0.000 [0.000, 0.354] \\
Program injected$^{\dagger}$ & 100-case positive-control & QMTester & Detected & 100/100 & coverage-by-constr. \\
Program injected & 100-case positive-control & MorphQ & Detected & 20/100 & 0.200 [0.133, 0.289] \\
Fixed variants & 7 Bugs4Q + 100 injected fixed & QMTester & False pos. & 0/107 & 0.000 [0.000, 0.035] \\
Correct programs$^{\flat}$ & ideal / noisy / calibrated sim & QMTester & False pos. & 0/50 \emph{each} & 0.000 [0.000, 0.071] \\
\bottomrule
\end{tabular}
\end{table*}


\subsection{RQ1: Real Program-Level Defects}
\label{sec:rq1}

\noindent\textbf{How many audited Bugs4Q program-level defects does QMTester detect?}

We \emph{included} a Bugs4Q subject only when a sound program-level relation could be justified (an input transform, a bijective canonicalization map, and a soundness argument tied to builder semantics). The denominators form one \emph{closed} chain: of the 42 Bugs4Q bugs~\cite{miao2023bugs4q}, 32 build and run under our pinned Qiskit (the diagnostic set, \S\ref{sec:diagnostic}); we audited \emph{all 32} (per-bug dispositions in the artifact), of which \emph{9} admit a sound builder-contract relation and 23 are out of scope (12 api-misuse, the rest output-formatting, platform/runtime, or wrong-gate; Table~\ref{tab:excluded-bugs4q}). Seven of the nine are runnable distribution-family subjects (Table~\ref{tab:audited-bugs4q}), all detected; the \texttt{program\_register\_shape} static check flags the eighth---\texttt{qiskit\_github\_12}, the \texttt{measure\_all} defect previously excluded for lacking a bijective map---and the ninth (a QPE/inverse-QFT case) awaits a sound relation. QMTester thus \emph{handles 8} of the 9 admitted defects (7/7 on the distribution families) and flags none of the matched fixed variants. Beyond Bugs4Q, mining recent Quantum Computing StackExchange (2022--2024) adds \emph{three} verified real defects (two QFT-round-trip, one input-permutation), each held to the \emph{same} bar---buggy detected, matched fix clean (\texttt{program\_mined\_manifest}). The real distribution base is thus \emph{ten}, all detected at $0$ false positives (Wilson $[0.722,1.000]$, up from $[0.646,1.000]$). The matched baselines extend: the reference-oracle detects all three (so QMTester matches the ceiling with \emph{no} oracle on all \emph{ten}), the static analyzer none ($0/10$); a circuit-level diagnostic incidentally fires on the QPE via swap-sensitivity (not a builder-contract detection, \S\ref{sec:diagnostic}).

\begin{table*}[t]
\centering
\caption{Audited Bugs4Q program-level subjects (in-scope). A subject is \emph{included} only when a sound program-level relation can be documented (source$\rightarrow$follow-up input transformation, bijective canonicalization map, and a soundness argument). For each, the buggy variant is detected and the fixed variant is clean. Full per-subject audit (source/follow-up inputs, canonical maps, soundness arguments, observed violations) is in the relation-soundness audit bundled with the artifact.}
\label{tab:audited-bugs4q}
\scriptsize
\setlength{\tabcolsep}{3pt}
\begin{tabular}{@{}p{0.225\textwidth}p{0.165\textwidth}p{0.35\textwidth}cc@{}}
\toprule
\textbf{Subject} & \textbf{Family (\texttt{program\_*})} & \textbf{Metamorphic transform (source\,$\rightarrow$\,follow-up input)} & \textbf{Buggy} & \textbf{Fixed} \\
\midrule
\texttt{github\_17\_meas\_order} & \texttt{classical\_remap} & Remap classical-register read order; canonicalize back to source layout. & detected & clean \\
\texttt{so\_5\_ccx\_role} & \texttt{input\_permutation} & Permute the two symmetric CCX control roles; invert on output. & detected & clean \\
\texttt{so\_1\_teleport\_feedback} & \texttt{ancilla\_uncompute} & Require clean teleportation feedback/ancilla state under declared uncompute. & detected & clean \\
\texttt{se\_3\_ccx\_uncompute} & \texttt{ancilla\_uncompute} & Complete the declared CCX ancilla uncompute; logical output invariant. & detected & clean \\
\texttt{se\_12\_meas\_endian} & \texttt{classical\_remap} & Flip classical-register endianness; canonicalize to source order. & detected & clean \\
\texttt{so\_7\_qft\_direction} & \texttt{qft\_round\_trip} & Compose QFT with its declared inverse; round trip is identity on input. & detected & clean \\
\texttt{se\_8\_grover\_uncompute} & \texttt{ancilla\_uncompute} & Require clean Grover ancilla after the declared uncompute. & detected & clean \\
\midrule
\multicolumn{3}{@{}l@{}}{\textbf{Included denominator: 7 \quad QMTester 7/7 detected \quad MorphQ 0/7 \quad fixed variants 0/7 false positives}} & & \\
\bottomrule
\end{tabular}
\end{table*}

\begin{table}[t]
\centering
\caption{Disposition of \emph{all} 32 build-and-run Bugs4Q bugs (full audit; per-bug detail in the artifact). Nine admit a sound builder-contract relation---eight now detected (Sec.~\ref{sec:rq1})---and 23 are out of scope by the categories below. With every bug dispositioned, the in-scope count is a population property, not an audit-coverage artifact.}
\label{tab:excluded-bugs4q}
\scriptsize
\setlength{\tabcolsep}{3pt}
\begin{tabular}{@{}p{0.62\columnwidth}c@{}}
\toprule
\textbf{Out-of-scope category} & \textbf{Count} \\
\midrule
API misuse (deprecated API, wrong import/attribute) & 12 \\
Output-formatting (visualization, global phase) & 4 \\
Platform/runtime (backend, transpiler, version) & 3 \\
Generic wrong-gate / missing-gate & 3 \\
Measurement is not a program contract for this bug & 1 \\
\midrule
\multicolumn{2}{@{}l@{}}{\textbf{Out of scope: 23 of 32} \quad (in scope: 9; detected: 8)} \\
\bottomrule
\end{tabular}
\end{table}


\noindent This is detection on the relation-applicable subset, not a suite-wide rate. Three matched baselines (Table~\ref{tab:program-baselines}): MorphQ is structurally blind ($0$); the \emph{deployable} static analyzer (LintQ-style register-shape check) also detects $0$---the circuits being equally well-formed, it flags only the one \emph{syntactic} \texttt{measure\_all} case (covered by \texttt{program\_register\_shape}); a reference-oracle that \emph{has} each known-correct version detects all, the ceiling but not deployable. QMTester reaches that ceiling with \emph{no} oracle; MorphQ's $0$ is circuit-level blindness---structural on the builder-only families, a catalog gap on the two circuit-expressible ones (\S\ref{sec:relations}). We are plain about what the statistics do \emph{not} carry: detection on all ten (deterministic) real subjects is carried by canonicalization (canonicalize-plus-raw-equality reaches $10/10$, \S\ref{sec:oracle-light}). The permutation test and Bonferroni earn their keep on \emph{specificity} and \texttt{parameter\_periodicity}, the one probabilistic family, which still has \emph{no real subject} (only the synthetic power curve, Table~\ref{tab:power-curve}).

\subsection{RQ2: Injected Positive Control}

\noindent\emph{Sanity check (coverage-by-construction).} Does each family fire on a balanced injected control, and what does QMTester's machinery add over a trivial rebuild-and-compare?

Because a mutant is admitted here only if detected (\S\ref{sec:threats}), the $100/100$ is coverage-by-construction (Table~\ref{tab:canonical-results}): a sanity check that each family fires on its own contract, \emph{not} a detection rate or contest. We retain it only as the within-method ablation substrate below. (MorphQ detects $20/100$---injected single-circuit \texttt{ancilla\_uncompute} mutants whose garbage \emph{is} circuit-visible---but \emph{none} of the three \emph{real} cross-boundary ancilla subjects, so its blindness holds for the real defects, not the injected proxy.)
\label{sec:oracle-light}

\noindent\emph{Oracle-light baselines.} To isolate QMTester's machinery, three \emph{program-level} baselines share the rebuild but each strip one component (Table~\ref{tab:program-baselines}): \emph{same input} (no transform), \emph{no canonicalization} (raw supports to the oracle), and \emph{raw equality}---same seven subjects and 100 mutants, identical shots/seed/correction.

\begin{table}[t]
\centering
\caption{Program-level baselines (detection on buggy/mutant, FP on matched fixed; 4096 shots, ideal sim). The reference-oracle differential is a non-deployable detection ceiling; the \emph{deployable} LintQ-style register-shape check catches only the \emph{syntactic} \texttt{measure\_all} case, none of the 7 \emph{semantic} defects. Indented rows are oracle-light ablations: canonicalize-plus-raw-equality alone already gives $7/7$, so the permutation test serves specificity, not detection.}
\label{tab:program-baselines}
\footnotesize
\setlength{\tabcolsep}{3pt}
\begin{tabular}{@{}p{0.585\columnwidth}cc@{}}
\toprule
Method & Det. & FP \\
\midrule
\multicolumn{3}{@{}l}{\emph{Audited Bugs4Q defects (7 buggy / 7 fixed)}}\\
MorphQ (circuit-level; blind) & 0/7 & --- \\
Static analyzer (\emph{deployable}; LintQ-style) & 0/7 & 0/7 \\
Reference-oracle (external; \emph{has} oracle) & 7/7 & 0/7 \\
QMTester (full; oracle-free) & \textbf{7/7} & \textbf{0/7} \\
\quad canonicalize $+$ raw equality (no test) & 7/7 & 0/7 \\
\quad same input $+\ \chi^2$ (no transform) & 0/7 & 0/7 \\
\quad no canonicalization $+\ \chi^2$ & 5/7 & 2/7 \\
\quad raw count-equality, no canon. & 5/7 & 2/7 \\
\midrule
\multicolumn{3}{@{}l}{\emph{Injected benchmark (100 mutant / 100 fixed)}}\\
QMTester (full) & \textbf{100/100} & \textbf{0/100} \\
\quad same input $+\ \chi^2$ & 0/100 & 0/100 \\
\quad no canonicalization $+\ \chi^2$ & 60/100 & 40/100 \\
\quad raw count-equality, no canon. & 60/100 & 40/100 \\
\bottomrule
\end{tabular}
\end{table}


The result isolates two contributions. First, the metamorphic transform is necessary: \emph{same input} detects $0/7$ and $0/100$. Second, canonicalization buys specificity: dropping it makes \emph{no canonicalization} and \emph{raw equality} false-positive on correct variants ($2/7$, $40/100$) \emph{and} miss buggy ones ($5/7$, $60/100$ detected), the un-undone representation change both masquerading as a fault and masking real divergences. The real-subject differential is exactly the two \texttt{classical\_remap} cases (2 points); the broader, inclusion-independent signal is the $40/100$ injected fixed-variant false positives QMTester's full pipeline holds at $0$---within-method evidence MorphQ cannot provide.

\subsection{RQ3: False Positives}

\noindent\textbf{What false-positive rate does QMTester report?}

Specificity is the only inclusion-\emph{independent} evidence in our evaluation, so we report it carefully. Across the 107 fixed program variants (7 Bugs4Q fixed subjects, 100 injected fixed variants) and the 50-program correct corpus, QMTester raises $0$ false positives under the ideal, noisy, and calibrated-noisy simulators (Table~\ref{tab:canonical-results}). The \emph{ideal} cell is degenerate (matched seeds make source and follow-up bit-identical, $p\equiv 1$); the informative cells are the noisy ones. Under the bundled depolarizing model ($p_1{=}10^{-3}$, $p_2{=}5{\times}10^{-3}$, $T_1{=}T_2{=}80\,\mu$s, $t_g{=}50$\,ns, density-matrix sim; \texttt{noise.json}) the smallest \emph{uncorrected} per-pair $p$ over all $3233$ pairs is $0.370$ (noisy) and $0.251$ (calibrated)---conservative, but not by a wide margin. To stress this we re-ran the corpus under a \emph{realistic} model adding $2\%$ symmetric readout error and a $1\%$ two-qubit crosstalk term: the smallest uncorrected $p$ falls to $0.0035$ (below the nominal $0.05$), yet source-level Bonferroni still yields $0/50$. So under realistic noise the per-pair margin is genuinely thin and the multiple-testing correction---not merely conservative nulls---is what holds specificity, a qualifier the four $0/50$ rows alone would hide. Concretely, the thin-margin subject (smallest realistic-noise $p=0.0035$) carries $m=96$ admitted pairs, so $\alpha/m=0.0005\le0.0035$; across the corpus $m\in[30,96]$, every per-subject threshold below its subject's smallest margin, so $0/50$ holds by correction, not by the thin-margin pair landing on a high-$m$ subject by luck.

\subsection{Statistical-Oracle Behavior}
\label{sec:oracle-behavior}

We verify two properties of the statistical oracle directly from the logged per-pair records. \emph{Empirical type-I:} across the 50 correct programs, none of the $3233$ per-pair null p-values falls below $0.05$ even uncorrected under any of the three simulators (minimum $1.000$/$0.370$/$0.251$ for ideal/noisy/calibrated-noisy). The nulls are thus not uniform but \emph{conservative}, the favorable direction for Remark~1, matching the observed $0/107$ and $0/50$ false-positive rates. \emph{Which branch decides:} the Monte-Carlo permutation branch, not the asymptotic $\chi^2$, fires for most program-level decisions ($85.7\%$ of Bugs4Q and $93.0\%$ of injected pairs), as the small quantum supports never satisfy the Cochran asymptotic condition (zero merges across all admitted pairs); we therefore describe the oracle as permutation-primary with an asymptotic $\chi^2$ fast path.

\subsection{RQ4: Sensitivity Floor (Detection vs.\ Effect Size)}

\noindent\textbf{At what effect size does detection start to fail?}

The families differ in output character. \texttt{input\_permutation}, \texttt{classical\_remap}, \texttt{ancilla\_uncompute}, and the real \texttt{qft\_round\_trip} subject all produce \emph{deterministic} basis-state divergence (TVD$=1.0$); the $0.74$--$0.94$ figure is the input-dependent effect size of the \emph{injected} QFT mutants, not the real subject. Only the probabilistic \texttt{parameter\_periodicity} family has a tunable effect. A $204$-mutant rotation-drift sweep (TVD$\in(0,0.06)$, \emph{not} filtered by detectability; true effect at 32{,}768 shots, detection at 4096) localizes the threshold via a fine sub-$0.03$ grid (Table~\ref{tab:power-curve}): the $50\%$/$80\%$ crossings fall in $[0.022,0.025)$/$[0.025,0.030)$. So the conservative nulls do not silently trade away sensitivity---the method detects reliably above a measured floor near TVD~$0.025$ (probabilistic family only).

\begin{table}[t]
\centering
\caption{Power curve (detection rate vs.\ effect size) for \texttt{program\_parameter\_periodicity} at 4096 shots, over a $204$-mutant rotation-drift sweep over TVD$\in(0,0.06)$, \emph{not} filtered by detectability (true effect at 32{,}768 shots). The fine sub-$0.03$ sweep \emph{localizes} the threshold: $50\%$ in $[0.022,0.025)$, $80\%$ in $[0.025,0.030)$. Other four families: (near-)deterministic (TVD$\approx 1$), no small-effect regime.}
\label{tab:power-curve}
\footnotesize
\begin{tabular}{lcc}
\toprule
Effect size (TVD) & Detected & Rate (Wilson 95\% CI) \\
\midrule
$[0.000,0.018)$ & $12/155$ & $0.08$~$[0.04,0.13]$ \\
$[0.018,0.022)$ & $5/12$  & $0.42$~$[0.19,0.68]$ \\
$[0.022,0.025)$ & $7/10$  & $0.70$~$[0.40,0.89]$ \\
$[0.025,0.060)$ & $27/27$ & $1.00$~$[0.88,1.00]$ \\
\bottomrule
\end{tabular}
\end{table}


\subsection{Diagnostic Circuit-Level Run}
\label{sec:diagnostic}

For transparency, a circuit-level-only configuration (identity/swap/phase/equivalence rewritings) shows \emph{no} advantage over MorphQ---$2/32$ on Bugs4Q, $17/250$ injected, identical for both, all from swap---motivating the program-level design (not a contribution claim).
